% !PRINTED 121
\weberpage{121}
%<<CONT
and the function $\sigma(u)$ itself may accordingly also be denoted by $\sigma_{11}(u)$.

For these functions there follow from (4) the characteristic period equations

\weq{6}{\begin{aligned}
\sigma_{g_1,g_2}(u+\omega_1) &= (-1)^{g_1}\,e^{\eta_1(2u+\omega_1)}\,\sigma_{g_1,g_2}(u),\\
\sigma_{g_1,g_2}(u+\omega_2) &= (-1)^{g_2}\,e^{\eta_2(2u+\omega_2)}\,\sigma_{g_1,g_2}(u),\\
\sigma_{g_1,g_2}(u+a\omega_1+b\omega_2) &= (-1)^{g_1 a+g_2 b+ab}\,e^{(a\eta_1+b\eta_2)(2u+a\omega_1+b\omega_2)}\,\sigma_{g_1,g_2}(u).
\end{aligned}}

By the application of a linear transformation the three functions $\sigma_{00}, \sigma_{01}, \sigma_{10}$ are permuted among one another, as formula (2) teaches [or \S\,35, (1)]. According to this permutation the linear transformations split into \gesperrt{six classes}, the first of which comprises all those transformations that leave the functions $\sigma_{00}, \sigma_{01}, \sigma_{10}$ unchanged. These are characterized by the fact that $\alpha, \delta$ are odd and $\beta, \gamma$ even numbers, which we write briefly thus:

\[
\begin{pmatrix}\alpha, & \beta\\ \gamma, & \delta\end{pmatrix} \equiv \begin{pmatrix}1, & 0\\ 0, & 1\end{pmatrix} (\mathrm{mod}\ 2).
\]

Accordingly the six classes of linear transformations are to be characterized as follows:

\weq{7}{\begin{aligned}
&\text{I.}\quad \begin{pmatrix}\alpha, & \beta\\ \gamma, & \delta\end{pmatrix} \equiv \begin{pmatrix}1, & 0\\ 0, & 1\end{pmatrix} (\mathrm{mod}\ 2)\quad (00, 01, 10),\\
&\text{II.}\quad \text{,,} \equiv \begin{pmatrix}0, & 1\\ -1, & 0\end{pmatrix} \quad \text{,,}\quad (00, 10, 01),\\
&\text{III.}\quad \text{,,} \equiv \begin{pmatrix}1, & 1\\ -1, & 0\end{pmatrix} \quad \text{,,}\quad (10, 00, 01),\\
&\text{IV.}\quad \text{,,} \equiv \begin{pmatrix}1, & 1\\ 0, & 1\end{pmatrix} \quad \text{,,}\quad (10, 01, 00),\\
&\text{V.}\quad \text{,,} \equiv \begin{pmatrix}1, & 0\\ 1, & 1\end{pmatrix} \quad \text{,,}\quad (01, 00, 10),\\
&\text{VI.}\quad \text{,,} \equiv \begin{pmatrix}0, & 1\\ -1, & 1\end{pmatrix} \quad \text{,,}\quad (01, 10, 00),
\end{aligned}}
where in the last column the permutation of the characteristics 00, 01, 10 occurring in each case is set out.

The transformations of the first class form a group $\mathfrak{A}$ contained in the group of \gesperrt{all} linear substitutions, and hence a subgroup of $\mathfrak{L}$ (\S\,28). If we set
%CONT>>
